% =========================================================================
% Capítulo 10 — Descripción de tareas, fases, equipos o procedimientos y
% diagrama de Gantt/cronograma (si procede)
% Metodología de trabajo (decisión SCRUM vs. tradicional/RUP, justificada)
% y las tareas/fases del desarrollo con su cronograma.
% OJO: el Gantt del TFG en sí (obligatorio) tiene su propio capítulo (13).
% =========================================================================

\chapter{Descripción de tareas, fases y procedimientos y diagrama de Gantt}
\label{ch:tareas}

Este capítulo justifica la metodología con la que se planifica y ejecuta el trabajo, define cómo se estima el esfuerzo y presenta el \eng{Product Backlog} inicial y las fases del desarrollo con su cronograma.
La representación gráfica del cronograma completo del \ac{TFG} (incluida la memoria y los hitos académicos) se recoge en el capítulo \ref{ch:gantt-tfg}.

\section{Metodología de trabajo}
\label{sec:metodologia}

\subsection{Restricciones que la metodología debe satisfacer}
\label{sec:metodologia-restricciones}

La elección no parte de preferencias abstractas, sino de las condiciones reales de este proyecto:

\begin{itemize}
  \item \textbf{Alcance emergente.}
        El modelo de datos y las entidades se diseñan de forma incremental y meditada durante el desarrollo, una decisión ya tomada del proyecto; una red completa de tareas definida por adelantado quedaría obsoleta en semanas.
  \item \textbf{Disponibilidad variable y conocida.}
        El verano está condicionado por compromisos laborales, septiembre por el traslado a una estancia Erasmus que se prolonga hasta finales de enero, y desde febrero el trabajo convive con las posibles prácticas de empresa; la metodología debe absorber esa variación —y el riesgo de que sea aún mayor (R8, capítulo \ref{ch:riesgos})— sin romperse.
  \item \textbf{Una única persona.}
        No hay roles que repartir ni reuniones que sincronizar; toda ceremonia debe justificar su coste ante un equipo de uno.
  \item \textbf{Dos entregables simultáneos.}
        Producto y memoria evolucionan a la vez, así que la redacción debe ser trabajo planificado de primera clase, no un resto.
  \item \textbf{Telemetría para esta memoria.}
        El método debe producir por sí mismo los datos de seguimiento (horas, velocidad, desviaciones) que alimentan este capítulo y el \ref{ch:gantt-tfg}.
  \item \textbf{Marco académico.}
        La dedicación de referencia del \ac{TFG}, fijada en el apartado \ref{sec:contexto-marco} y repartida a lo largo de un curso completo, con la defensa prevista en la convocatoria de junio--julio de \thesisYear{} (capítulo \ref{ch:objetivos}).
\end{itemize}

\subsection{Alternativas consideradas y descartadas}
\label{sec:metodologia-alternativas}

\paragraph{Metodologías predictivas (cascada, \acs{RUP}) con planificación por precedencias.}
El enfoque tradicional de la gestión de proyectos (ciclos tipo \ac{RUP} y técnicas de planificación temporal como \ac{PERT}, Roy o los diagramas de precedencias) exige descomponer el proyecto completo en tareas con duraciones y dependencias estables antes de empezar \parencite{kruchten-rup, pmi-pmbok}.
Esa premisa falla aquí dos veces: el alcance es emergente por decisión de diseño, y la capacidad disponible varía a lo largo del periodo, de modo que la red de tareas exigiría re-planificación continua, un coste permanente sin valor añadido.
A ello se suma que \ac{RUP} prescribe una disciplina de artefactos y roles pensada para equipos, desproporcionada para una persona.
De este enfoque se rescata, no obstante, una pieza: un plan macro de fases e hitos con fechas (capítulo \ref{ch:gantt-tfg}), que da la visión de conjunto que una pila de historias no ofrece.

\paragraph{Scrum puro.}
Scrum es la metodología que el autor ya empleó, con resultados satisfactorios, en un proyecto de software desarrollado en equipo durante el grado, y por ello fue la primera candidata.
Sin embargo, Scrum se define explícitamente sobre un equipo con roles diferenciados (\eng{Product Owner}, \eng{Scrum Master}, \eng{Developers}) y ceremonias cuyo propósito es sincronizar y proteger a ese equipo \parencite{scrum-guide}.
En solitario, la mitad del marco queda vacío: no hay a quién hacer de contrapeso en la negociación del alcance, la reunión diaria de sincronización (\eng{daily}) no tiene interlocutor y la demostración del avance al cierre de cada iteración (\eng{review}) no tiene público.
Además, el compromiso cerrado de cada \eng{sprint} encaja mal con una capacidad que se sabe variable: un mes de baja dedicación encadenaría «fracasos» de \eng{sprint} que no informan de nada.
Declarar «se usó Scrum» sería, en rigor, falso; lo honesto es declarar qué elementos de Scrum se adoptan y por qué.

\paragraph{Kanban puro.}
El método Kanban (flujo continuo, trabajo visualizado, límite de trabajo en curso) encaja de forma natural con la disponibilidad variable: si una semana rinde menos, fluye menos trabajo, sin ningún compromiso roto \parencite{anderson-kanban}.
Se descarta como método único porque carece de cadencia: sin iteraciones no hay ritmo fijo de re-planificación ni retrospectiva, y no se genera la medida periódica de velocidad que este \ac{TFG} necesita como telemetría para la memoria y para corregir su primera experiencia de estimación en solitario (riesgo R2).

\subsection{Decisión: método híbrido (Scrumban) adaptado a una persona}
\label{sec:metodologia-decision}

Se adopta un híbrido de la familia Scrumban \parencite{ladas-scrumban, kniberg-scrumban}: la estructura y la telemetría de Scrum que sobreviven al contexto unipersonal, sobre una mecánica de flujo Kanban que absorbe la variabilidad.
La tabla \ref{tab:metodologia-elementos} resume qué se toma de cada marco, con qué adaptación y qué añade el método por su cuenta.

\begin{table}[htbp]
  \centering
  \small
  % Título encima de la tabla, como en las longtable: posición unificada en
  % toda la memoria (revisión del 24-08-2026); igual en el resto de tablas.
  \caption{Elementos del método híbrido y su adaptación al contexto unipersonal.}
  \label{tab:metodologia-elementos}
  \begin{tabular}{p{3.4cm}p{1.6cm}p{8.2cm}}
    \toprule
    \textbf{Elemento} & \textbf{Origen} & \textbf{Adaptación a una persona} \\
    \midrule
    \eng{Product Backlog} & Scrum & Historias de usuario agrupadas en épicas clasificadas por prioridad en cuatro niveles (imprescindible, deseable, opcional, excluida; sección \ref{sec:backlog}); toda idea nueva entra por aquí (control del alcance, R1). \\
    Iteraciones & Scrum & \eng{Sprints} semanales (lunes--domingo) con capacidad variable declarada por sprint. \\
    Naturaleza del compromiso & Kanban & El compromiso es el orden del \eng{backlog} y el límite de \ac{WIP}, no un contrato cerrado de \eng{sprint}; el contenido semanal es una proyección. \\
    Límite de \acs{WIP} & Kanban & Máximo 2 elementos en curso: una historia de producto y, como mucho, la redacción de la memoria; nunca dos historias de producto a medias. \\
    Planificación de \eng{sprint} & Scrum & Sesión individual de unos 20 minutos cada lunes: fijar capacidad, seleccionar del \eng{backlog} y anotar la lista de trabajo de la semana en la entrada del diario de ese día. \\
    \eng{Daily} & Scrum & Sustituida por el procedimiento ya vigente de abrir cada sesión releyendo la última entrada del diario de desarrollo. \\
    \eng{Review} & Scrum & Verificación del incremento contra los criterios de cierre de cada historia (sección \ref{sec:dod}) y, cada quince días, entrega de la versión vigente al tutor (sección \ref{sec:metodologia-tutor}). \\
    Retrospectiva & Scrum & Sección fija de la entrada del diario al cierre del \eng{sprint}, que incluye la revisión de riesgos (capítulo \ref{ch:riesgos}). \\
    Revisión de coherencia & Propio & Al cierre de cada \eng{sprint}, revisión automática de la coherencia entre la documentación, la memoria y el código; el autor decide el lunes qué se corrige (sección \ref{sec:metodologia-coherencia}). \\
    Velocidad & Scrum & Puntos completados por \eng{sprint}; recalibra la estimación de forma continua. \\
    \bottomrule
  \end{tabular}
\end{table}

La duración de una semana por \eng{sprint} (frente a las dos semanas empleadas en aquella experiencia previa en equipo) se justifica por tres razones.
Primera: la realimentación de velocidad llega al doble de frecuencia justo cuando más falta hace, en las primeras semanas de la primera estimación en solitario del autor.
Segunda: el grano semanal permite modular la capacidad de cada periodo (verano con trabajo, estancia Erasmus, posibles prácticas de empresa) sin tocar el método: cambia la capacidad, no el ritmo.
Tercera: el argumento clásico contra los \eng{sprints} cortos, el coste proporcional de las ceremonias, desaparece en solitario, donde todas las ceremonias juntas cuestan menos de una hora semanal.

\subsection{Herramienta de gestión y su evolución}
\label{sec:metodologia-herramienta}

El método arrancó con un tablero Kanban en GitHub Projects.
Cada historia era una ficha de trabajo de GitHub (\eng{issue}), con sus puntos y su épica de pertenencia (el bloque de funcionalidad que la agrupa; sección \ref{sec:backlog}) como etiquetas, y cada fase era un hito (\eng{milestone}).
El tablero tenía cuatro columnas (\eng{Backlog}, Preparado, En curso con límite de \acs{WIP} 2, y Hecho), y cada cambio de código cerraba la ficha de su historia, enlazando historia y código.

Ese tablero se usó durante toda la fase F1 y se abandonó al retomar el trabajo tras la pausa del traslado (03/10/2026), junto con las fichas y los hitos, por tres motivos que la propia fase dejó a la vista.
Primero, era innecesario: en un proyecto de una sola persona duplicaba lo que ya registraban el diario de desarrollo y este capítulo, de modo que no aportaba información nueva.
Segundo, metía ruido en la planificación: el trabajo real de cada semana quedaba repartido entre el tablero, el diario y la memoria, y mover fichas no informaba de nada que el diario no dijera ya.
Tercero, suponía una carga de trabajo extra sin retorno, justo en el recurso que la fase demostró más escaso: las horas disponibles (capítulo \ref{ch:gantt-tfg}).

Desde entonces, el \eng{backlog} vive en este capítulo (tablas \ref{tab:backlog-epicas} y \ref{tab:backlog-f1}), la planificación de cada lunes se anota como lista de trabajo de la semana en la entrada del diario de ese día, y el límite de una sola historia de producto en curso se mantiene como regla, ya sin tablero que lo visualice.
La trazabilidad entre historias y código se conserva por otra vía: cada cambio de código nombra la historia que implementa, y el diario la registra.
Las fichas y los hitos de F1 se conservan cerrados en el repositorio, como evidencia de esa primera etapa del método.
Esta memoria conserva las instantáneas relevantes de la planificación: el \eng{backlog} inicial aquí, y el estado final y las desviaciones en los capítulos \ref{ch:resultados} y \ref{ch:gantt-tfg}.

\subsection{Revisión periódica con el tutor}
\label{sec:metodologia-tutor}

Al retomar el trabajo, el método incorpora una revisión externa periódica: los días 1 y 15 de cada mes, el autor envía al tutor la versión vigente de la memoria con un resumen de los cambios y novedades desde el envío anterior.
La primera entrega coincidió con la reanudación del trabajo, el 05/10/2026, y la segunda se fijó para el 15/10/2026, de modo que la cadencia quincenal arranca ese mismo mes.
La cadencia cumple dos funciones.
La primera es dar a la revisión del avance el destinatario que en solitario le faltaba (sección \ref{sec:metodologia-alternativas}): alguien externo que valora el progreso y propone mejoras mientras aún pueden incorporarse.
La segunda es ejercer una disciplina indirecta: un compromiso fijo cada dos semanas sostiene un ritmo de trabajo constante, que es precisamente lo que faltó en la fase F1.

\subsection{Revisión semanal de la coherencia documental}
\label{sec:metodologia-coherencia}

Un proyecto documentado con este nivel de detalle acumula muchas fuentes que cuentan lo mismo desde ángulos distintos: las reglas de trabajo persistentes, el diario de desarrollo, la documentación del repositorio, los diagramas, los documentos de riesgos, los comentarios del código y esta misma memoria.
Con cada decisión nueva alguna de ellas se queda atrás, y el desfase no se nota hasta que alguien lo encuentra por casualidad: una versión que ya no es la instalada, una regla que un documento da por vigente y otro por superada, una cifra que no coincide entre dos capítulos.
Por eso, desde el cierre del séptimo \eng{sprint}, el final de cada \eng{sprint} incluye una revisión automática de la coherencia entre todas esas fuentes y con el código real.

La revisión la hace un procedimiento asistido por \ac{IA}, preparado como un comando propio que se lanza el domingo por la noche y trabaja sin supervisión.
Primero reúne los hechos objetivos del repositorio con pequeños programas que no usan \ac{IA}: versiones instaladas, estado de la integración continua y del despliegue, fechas y \eng{sprints} calculados, y citas y acrónimos de la memoria.
Después reparte la comparación entre seis agentes que trabajan en paralelo, cada uno dedicado a una dimensión: tecnologías y modelo de datos; arquitectura y documentación del repositorio; producto y método; riesgos y presupuesto; la propia memoria; y el diario de desarrollo junto con la privacidad y los comentarios del código.
Cada contradicción grave pasa además por un agente que intenta refutarla releyendo las fuentes, para no dar por buenos los falsos positivos.
El resultado es un informe fechado con cada contradicción, la fuente que prevalece según unas reglas de precedencia escritas y la corrección propuesta.
El procedimiento no modifica nada: en la planificación del lunes siguiente, el autor decide qué correcciones se aplican y resuelve las que exigen una decisión de contenido.

La escala del procedimiento se acotó a propósito.
Un primer diseño, con quince dimensiones y un agente verificador por cada hallazgo, podía movilizar varios cientos de agentes en cada ejecución: un coste razonable para una auditoría puntual, pero insostenible para una práctica semanal sobre un plan de suscripción con límites de uso.
Se rehízo con seis dimensiones, una única verificación agrupada de los hallazgos graves y modelos más ligeros en las comprobaciones mecánicas, de modo que cada ejecución emplea menos de diez agentes.

La práctica refuerza la mitigación del riesgo R3, la memoria relegada frente al desarrollo (capítulo \ref{ch:riesgos}): sustituye las puestas al día concentradas, como la que exigió la reanudación tras la primera pausa, por una comprobación continua de coste acotado.

\section{Estimación del esfuerzo}
\label{sec:estimacion}

El esfuerzo se estima en puntos de historia (\eng{story points}), una medida relativa de tamaño, con la escala de Fibonacci modificada (1, 2, 3, 5, 8 y 13).
La estimación relativa se elige frente a la estimación directa en horas porque comparar es más fiable que predecir duraciones absolutas, y porque desacopla el tamaño del trabajo de la capacidad disponible, que aquí varía por diseño \parencite{cohn-estimating}.
Al no existir equipo, no procede la estimación por votación simultánea del equipo (\eng{planning poker}); en su lugar, cada historia se estima por triangulación contra historias de referencia ya conocidas.

La historia es además la única unidad que recibe puntos, una decisión que conviene justificar porque se aparta a sabiendas de la práctica que el autor conoció en el proyecto en equipo ya mencionado (sección \ref{sec:metodologia-alternativas}), donde cada historia se desglosaba en tareas y cada tarea recibía su propia estimación en puntos.
Ese segundo nivel de estimación cumple en un equipo una función concreta: repartir el trabajo de una misma historia entre varias personas, equilibrar la carga de cada miembro y hacer visible el avance individual dentro de la iteración.
En un proyecto de una sola persona esa función desaparece por completo —no hay trabajo que repartir ni cargas que equilibrar—, y de la doble estimación solo quedarían los costes: un trámite añadido en cada historia y, sobre todo, el riesgo de contar el mismo trabajo dos veces y distorsionar la velocidad, la única telemetría con la que este plan se recalibra.
Las tareas, en cambio, sí se conservan cuando ayudan: al ejecutar una historia, sus subtareas técnicas se anotan como lista de comprobación en la planificación semanal del diario de desarrollo (durante F1, dentro de la ficha de la historia en el tablero; sección \ref{sec:metodologia-herramienta}), como guía de ejecución y rastro del avance, pero sin puntos propios.

Como hipótesis inicial de equivalencia se toman unas 2,7 horas por punto, de modo que la velocidad esperada de cada \eng{sprint} se deriva de su capacidad declarada (por ejemplo, unos 4 puntos en un \eng{sprint} de verano de 10 horas).
Esta equivalencia sirve únicamente para el encaje macro del plan; a partir del tercer \eng{sprint} la manda la velocidad real medida, y la hipótesis se recalibra en cada retrospectiva (riesgo R2).
La primera recalibración con tres \eng{sprints} medidos (cierre de S3, 09/08/2026) la fija en 2,0 horas por punto: el esfuerzo observado rondó las 1,2 —la asistencia de \ac{IA} aceleró la ejecución muy por encima de lo previsto—, y se adopta un valor deliberadamente prudente porque la muestra es solo de infraestructura; las épicas de mapa, tiempo real e \ac{IA} no tienen por qué rendir igual.

\section{Criterios de cierre de una historia}
\label{sec:dod}

Una historia solo se considera terminada cuando cumple todas las condiciones siguientes (lo que Scrum denomina la \emph{Definición de Hecho} \parencite{scrum-guide}):

\begin{itemize}
  \item El análisis estático y la comprobación de tipos del \eng{monorepo} pasan en verde.
  \item Las pruebas de la historia existen y pasan, y la suite completa no sufre regresiones.
  \item Si la historia toca autenticación, autorización o datos personales, sus pruebas incluyen los casos de abuso correspondientes (entradas inválidas o maliciosas, accesos no autorizados): el objetivo de seguridad del capítulo \ref{ch:objetivos} se verifica historia a historia, no al final (riesgo R14).
  \item El cambio está integrado en la rama principal.
  \item Las decisiones tomadas quedan anotadas en el diario de desarrollo, con la alternativa descartada y su porqué.
  \item Si la historia afecta a un capítulo de esta memoria, el capítulo queda actualizado (mitigación del riesgo R3).
\end{itemize}

Esta lista es el criterio con el que se verifica cada incremento —la parte del producto terminada en el \eng{sprint}— y protege los objetivos de calidad y trazabilidad del capítulo \ref{ch:objetivos}.

\section{\eng{Product Backlog} inicial}
\label{sec:backlog}

El \eng{backlog} inicial se construye descomponiendo las reglas de producto de \bukit en \emph{épicas} —bloques grandes de funcionalidad que agrupan varias historias de usuario relacionadas y que son demasiado extensos para abordarse de una vez—, y estas en historias de usuario con el formato «como \emph{rol}, quiero \emph{acción}, para \emph{beneficio}» \parencite{cohn-user-stories}.
Cada épica se clasifica por prioridad en cuatro niveles, siguiendo la técnica conocida como MoSCoW \parencite{clegg-moscow}: \emph{imprescindible} para el alcance comprometido del que dependen los criterios de éxito, \emph{deseable} para lo incluido en el alcance sin criterio de éxito asociado, \emph{opcional} para los objetivos declarados secundarios, y \emph{excluida} para lo explícitamente fuera del alcance (sección \ref{sec:alcance-fuera}), que no se repite aquí.
Esta clasificación traduce el alcance del capítulo \ref{ch:objetivos} a una regla operativa simple: lo imprescindible se protege siempre, y lo deseable y lo opcional solo entran si sobra velocidad.

\begin{table}[htbp]
  \centering
  \small
  \caption{Épicas del \eng{Product Backlog} inicial, con estimación gruesa y fase prevista.}
  \label{tab:backlog-epicas}
  \begin{tabular}{clccc}
    \toprule
    \textbf{Id} & \textbf{Épica} & \textbf{Prioridad} & \textbf{Puntos} & \textbf{Fase} \\
    \midrule
    EP01 & Fundaciones técnicas e integración continua      & Imprescindible & 8  & F1 \\
    EP02 & Cuentas y autenticación                          & Imprescindible & 13 & F1 \\
    EP03 & Andamiaje del cliente web y sesión               & Imprescindible & 8  & F2 \\
    EP04 & Listas y metas                                   & Imprescindible & 13 & F2 \\
    EP05 & Publicaciones ligadas a listas                   & Imprescindible & 13 & F2 \\
    EP06 & Grafo social y \eng{feed} de seguidos            & Imprescindible & 10 & F3 \\
    EP07 & Perfil y verificación de identidad               & Imprescindible & 8  & F3 \\
    EP08 & \eng{Feed} personalizado y buscador              & Imprescindible & 13 & F4 \\
    EP09 & Mapa social                                      & Imprescindible & 10 & F3 \\
    EP10 & Afinidad entre personas usuarias                 & Imprescindible & 8  & F3 \\
    EP11 & Mensajería y notificaciones                      & Imprescindible & 13 & F3 \\
    EP12 & Pruebas \acs{E2E} y despliegue continuo          & Imprescindible & 8  & F4 \\
    EP13 & \acs{IA} de estimación de plazos                 & Opcional       & 8  & --- \\
    EP14 & Aplicación móvil parcial                         & Opcional       & 13 & --- \\
    EP15 & Memoria y gestión del proyecto (transversal)     & Imprescindible & 21 & Todas \\
    EP16 & Subida y reproducción de vídeo (servicio gestionado) & Opcional   & 8  & --- \\
    EP17 & Publicaciones patrocinadas simuladas             & Opcional       & 5  & --- \\
    \bottomrule
  \end{tabular}
\end{table}

Las estimaciones de la tabla \ref{tab:backlog-epicas} son deliberadamente gruesas y no se someten a la escala de historias de la sección \ref{sec:estimacion}: anticipan la suma de los puntos de las historias que cada épica contendrá, y solo estas últimas se estiman con la escala declarada.
Siguiendo la elaboración progresiva propia de los métodos ágiles, solo se detallan en historias las épicas de la fase inminente, y el resto se refina justo antes de abordarse, con el conocimiento acumulado hasta entonces.
La tabla \ref{tab:backlog-f1} recoge el desglose completo de las épicas de fundaciones (EP01--EP03), ya refinado: EP01 y EP02 forman la fase F1, y EP03 abre la F2.

\begin{table}[htbp]
  \centering
  \small
  \caption{Desglose en historias de las épicas de fundaciones EP01--EP03, 29 puntos.}
  \label{tab:backlog-f1}
  \begin{tabular}{clcc}
    \toprule
    \textbf{Id} & \textbf{Historia (resumen)} & \textbf{Épica} & \textbf{Puntos} \\
    \midrule
    US-01 & Cadena de integración continua: lint, tipos y pruebas     & EP01 & 3 \\
    US-02 & Manejo centralizado de errores y registro de sucesos      & EP01 & 3 \\
    US-03 & Despliegue del esqueleto de la \acs{API}                  & EP01 & 2 \\
    US-04 & Modelo de datos de usuario y registro de cuentas          & EP02 & 3 \\
    US-05 & Inicio de sesión con \acs{JWT} de acceso y refresco       & EP02 & 3 \\
    US-06 & Rotación y revocación de los \eng{tokens} de refresco     & EP02 & 5 \\
    US-07 & Doble capa de limitación de peticiones                    & EP02 & 2 \\
    US-08 & Andamiaje web con enrutado de las cuatro pestañas         & EP03 & 3 \\
    US-09 & Cliente \acs{API} tipado y gestión de sesión en la web    & EP03 & 3 \\
    US-10 & Pantallas de registro e inicio de sesión                  & EP03 & 2 \\
    \bottomrule
  \end{tabular}
\end{table}

En la tabla \ref{tab:backlog-f1} cada historia aparece resumida en una línea, pero antes de entrar en un \eng{sprint} toda historia se desarrolla hasta un nivel de detalle que la hace verificable.
Como muestra de ese nivel se desarrolla completa la que sería la primera historia de EP04, continuando la numeración de la tabla:

\begin{quote}
  \textbf{US-11 (EP04, 5 puntos).}
  Como persona usuaria registrada, quiero crear una lista de metas, cada una con su fecha límite, para convertir mis deseos en compromisos medibles.
  \emph{Criterios de aceptación:} (1) una cuenta gratuita no puede superar las tres listas en curso y el exceso se rechaza con un error explícito; (2) la validación vive una sola vez en el paquete compartido y la usan tanto la \ac{API} como el cliente web; (3) toda lista exige al menos una meta y cada meta una fecha límite; (4) el flujo queda cubierto por pruebas de integración.
\end{quote}

Este desarrollo es una anticipación puramente ilustrativa, no el refinamiento de EP04, que se hará en su momento como el de cualquier otra épica.
Sus 5 puntos tampoco añaden trabajo nuevo al inventario: forman parte de los 13 estimados en grueso para EP04 en la tabla \ref{tab:backlog-epicas}.

\section{Fases y cronograma del desarrollo}
\label{sec:fases}

El periodo de ejecución se organiza en cuatro fases que agrupan 37 \eng{sprints} semanales (lunes--domingo), del 20/07/2026 al 16/05/\thesisYear, precedidas por una fase de planificación (F0, con un único \eng{sprint} S0, del 14 al 19 de julio de 2026) en la que se elaboran los objetivos, el plan de riesgos y este plan: la planificación inicial que la metodología elegida toma prestada del enfoque tradicional.
El calendario incorpora además dos pausas planificadas sin \eng{sprints}, correspondientes a los traslados de ida y vuelta de la estancia Erasmus: planificar semanas de capacidad cero no aporta información y solo introduciría ruido en la medida de velocidad.
Cada fase cierra con un incremento demostrable, que actúa como hito de control.
Estas fechas —incluida la de cierre del desarrollo, el 16/05/\thesisYear— son una base orientativa del plan, no un compromiso inamovible: fijan el encaje macro y los hitos de control, y cualquier re-base que las circunstancias exijan queda documentada con su motivo en este capítulo y reflejada en los capítulos \ref{ch:objetivos} y \ref{ch:gantt-tfg}.

Durante la ejecución, el calendario se re-basó una vez, el 03/10/2026.
F1 cerró en su fecha prevista (30/08/2026) con 14 de sus 21 puntos entregados, y las dos historias de EP02 que quedaron pendientes (US-06 y US-07) pasaron al inicio de F2.
La primera pausa, prevista hasta el 20/09/2026, se amplió hasta el 04/10/2026 porque el traslado y la instalación en la estancia Erasmus se alargaron más de lo previsto.
En consecuencia, F2 arrancó el 05/10/2026 con dieciséis \eng{sprints} en lugar de dieciocho (80 horas en lugar de 90), la numeración de los \eng{sprints} posteriores se desplazó dos posiciones y la fecha de cierre del desarrollo se mantuvo.
El balance de F1 frente a lo previsto se recoge en el capítulo \ref{ch:gantt-tfg}.

\begin{table}[htbp]
  \centering
  \small
  \setlength{\tabcolsep}{4pt}
  \caption{Fases del desarrollo con su capacidad estimada (la épica transversal EP15 se reparte entre todas).}
  \label{tab:fases}
  \begin{tabular}{clllc}
    \toprule
    \textbf{Fase} & \textbf{Contenido} & \textbf{Sprints} & \textbf{Fechas} & \textbf{Horas} \\
    \midrule
    F0  & Planificación: objetivos, riesgos, metodología    & S0       & 14/07 -- 19/07/2026 & $\sim$20 \\
    F1  & Fundaciones de la \acs{API}: EP01, EP02           & S1--S6   & 20/07 -- 30/08/2026 & 48 \\
    --- & Pausa: traslado e instalación en la estancia Erasmus & --- & 31/08 -- 04/10/2026 & 0 \\
    F2  & Cliente web, listas y publicaciones: EP03--EP05   & S7--S22  & 05/10/2026 -- 24/01/\thesisYear & 80 \\
    --- & Pausa: regreso e instalación en Bilbao            & ---      & 25/01 -- 31/01/\thesisYear & 0 \\
    F3  & Núcleo social y diferenciadores: EP06, EP07, EP09--EP11 & S23--S30 & 01/02 -- 28/03/\thesisYear & 120 \\
    F4  & Cierre de producto y memoria: EP08, EP12          & S31--S37 & 29/03 -- 16/05/\thesisYear & 105 \\
    \bottomrule
  \end{tabular}
\end{table}

La capacidad declarada refleja la disponibilidad real esperada de cada periodo: el verano se planificó con 10 horas semanales condicionadas por los compromisos laborales, y la realidad de tres \eng{sprints} la re-basó a 6 desde S4 (09/08/2026, contingencia del riesgo R8: la carga laboral resultó estructural, no puntual); unas 5 horas semanales de promedio durante la estancia Erasmus, donde las semanas de viaje se declaran a cero en su planificación de lunes y la mecánica de flujo las absorbe sin ningún compromiso roto; y unas 15 horas semanales desde febrero, compatibles con las posibles prácticas de empresa.
Que la capacidad real confirme estas hipótesis es exactamente lo que vigila el riesgo R8, con las horas registradas en el diario como evidencia.
Siguiendo la misma elaboración progresiva aplicada al \eng{backlog}, solo se detallan los \eng{sprints} de la fase inminente; los de fases posteriores se concretan en su planificación de lunes con el conocimiento acumulado hasta entonces.

\begin{table}[htbp]
  \centering
  \small
  \caption{Calendario de \eng{sprints} de F0 y F1 con su capacidad; el tema es una proyección, el compromiso real es el orden del \eng{backlog}.}
  \label{tab:sprints}
  \begin{tabular}{cllc}
    \toprule
    \textbf{Sprint} & \textbf{Fechas} & \textbf{Tema orientativo} & \textbf{Horas} \\
    \midrule
    S0 & 14/07 -- 19/07 & Planificación: objetivos, riesgos, metodología y este plan & $\sim$20 \\
    S1 & 20/07 -- 26/07 & Integración continua mínima (EP01)                         & 10 \\
    S2 & 27/07 -- 02/08 & Esqueleto de la \acs{API}, errores y registro de sucesos (EP01) & 10 \\
    S3 & 03/08 -- 09/08 & Modelo de usuario y registro de cuentas (EP02)             & 10 \\
    S4 & 10/08 -- 16/08 & Inicio de sesión con \eng{tokens} de acceso y refresco (EP02) & 6 \\
    S5 & 17/08 -- 23/08 & Rotación y revocación de \eng{tokens} (EP02)               & 6 \\
    S6 & 24/08 -- 30/08 & Limitación de peticiones; \emph{cierre de F1}              & 6 \\
    \bottomrule
  \end{tabular}
\end{table}

\subsection{Encaje del plan y margen declarado}
\label{sec:fases-tension}

El plan inicial encajaba al límite: el inventario comprometido (prioridad imprescindible) sumaba 138 puntos, y la capacidad nominal original (375 horas, a la hipótesis inicial de 2,7 horas por punto) daba unos 139 — sin holgura interna, con F2 ligeramente sobrecargada y F4 con excedente deliberado, reservado para el cierre de la memoria y la parte final de la épica transversal EP15.
El inventario comprometido vigente suma 146 puntos: los 138 iniciales más EP10 (afinidad), reclasificada de deseable a imprescindible el 24/08/2026 con fase asignada (F3, junto al mapa que la consume) porque el criterio de éxito del mapa de descubrimiento (capítulo \ref{ch:objetivos}) descansa en el cálculo de afinidad de esa épica, y mantenerla fuera del compromiso dejaba un criterio de éxito comprometido colgando de una épica no comprometida.
Ese encaje al límite fue la primera observación registrada del riesgo R2, y la primera recalibración con datos reales (cierre de S3) lo transformó: con la equivalencia en 2,0 horas por punto y la capacidad re-basada (353 horas tras las re-bases de F1 y de la primera pausa), los 146 puntos comprometidos suponen unas 292 horas — una holgura nominal en torno al 17\,\%.
Esa cifra corresponde al plan completo; tras el cierre de F1, que consumió sus 48 horas entregando 14 de sus 21 puntos, la holgura útil se mide sobre lo que queda: los 132 puntos pendientes suponen unas 264 horas frente a las 305 de capacidad de F2 a F4, es decir, una holgura en torno al 13\,\%.
La holgura se declara con la misma cautela con la que antes se declaró su ausencia: nace de una equivalencia medida solo sobre historias de infraestructura muy asistidas por \ac{IA}, y las épicas de mapa, tiempo real e \ac{IA} pueden devolverla a la realidad.
La decisión consciente sigue siendo triple: se mantiene el 16/05/\thesisYear{} como fecha objetivo de cierre del desarrollo, el margen real es el tramo desde esa fecha hasta el depósito de la convocatoria de junio--julio de \thesisYear, con la de septiembre como reserva (cronograma completo en el capítulo \ref{ch:gantt-tfg}), y las prioridades deseables y opcionales solo entran si la velocidad real consolida esa holgura.

Estas hipótesis de capacidad son deliberadamente conservadoras: describen el escenario razonable de \emph{menor} disponibilidad, y sus consecuencias —las fechas como techo del escenario prudente, no como suelo— se analizan en el capítulo \ref{ch:gantt-tfg}.
